\documentclass[../main]{subfiles}
\begin{document}
\section{Anexo:Sistemas Termodinámicos}
\subsection{Termodinámica Definición}
La Termodinámica es la rama de la Física que estudia a nivel macroscópico las transformaciones de la energía, y cómo esta energía puede convertirse en trabajo (movimiento). Históricamente, la Termodinámica nació en el siglo XIX de la necesidad de mejorar el rendimiento de las primeras máquinas térmicas fabricadas por el hombre durante la Revolución Industrial.

La Termodinámica clásica (que es la que se tratará en estas páginas) se desarrolló antes de que la estructura atómica fuera descubierta (a finales del siglo XIX), por lo que los resultados que arroja y los principios que trata son independientes de la estructura atómica y molecular de la materia.

El punto de partida de la mayor parte de consideraciones termodinámicas son las llamadas leyes o principios de la Termodinámica. En términos sencillos, estas leyes definen cómo tienen lugar las transformaciones de energía. Con el tiempo, han llegado a ser de las leyes más importantes de la ciencia.

Antes de entrar en el estudio de los principios de la termodinámica, es necesario introducir algunas nociones preliminares, como qué es un sistema termodinámico, cómo se describe, qué tipo de transformaciones puede experimentar, etc. Estos conceptos están resumidos en el siguiente cuadro:

\begin{center}
\includegraphics{images/02/intro_termo.png}    
\end{center}

\subsection{Sistema Termodinámico}

Un sistema termodinámico (también denominado sustancia de trabajo) se define como la parte del universo objeto de estudio. Un sistema termodinámico puede ser una célula, una persona, el vapor de una máquina de vapor, la mezcla de gasolina y aire en un motor térmico, la atmósfera terrestre, etc.

El sistema termodinámico puede estar separado del resto del universo (denominado alrededores del sistema) por paredes reales o imaginarias. En este último caso, el sistema objeto de estudio sería, por ejemplo, una parte de un sistema más grande. Las paredes que separan un sistema de sus alrededores pueden ser aislantes (llamadas paredes adiabáticas) o permitir el flujo de calor (diatérmicas).

Los sistemas termodinámicos pueden ser aislados, cerrados o abiertos.

Sistema aislado: es aquél que no intercambia ni materia ni energía con los alrededores.
Sistema cerrado: es aquél que intercambia energía (calor y trabajo) pero no materia con los alrededores (su masa permanece constante).
Sistema abierto: es aquél que intercambia energía y materia con los alrededores.
En la siguiente figura se han representado los distintos tipos de sistemas termodinámicos.

\begin{center}
\includegraphics{images/02/sistema_termo.png}    
\end{center}

A lo largo de estas páginas trataremos los sistemas cerrados.

Cuando un sistema está aislado y se le deja evolucionar un tiempo suficiente, se observa que las variables termodinámicas que describen su estado no varían. La temperatura en todos los puntos del sistema es la misma, así como la presión. En esta situación se dice que el sistema está en equilibrio termodinámico.

\subsection{Equilibrio termodinámico}

En Termodinámica se dice que un sistema se encuentra en equilibrio termodinámico cuando las variables intensivas que describen su estado no varían a lo largo del tiempo.

Cuando un sistema no está aislado, el equilibrio termodinámico se define en relación con los alrededores del sistema. Para que un sistema esté en equilibrio, los valores de las variables que describen su estado deben tomar el mismo valor para el sistema y para sus alrededores. Cuando un sistema cerrado está en equilibrio, debe estar simultáneamente en equilibrio térmico y mecánico.

Equilibrio térmico: la temperatura del sistema es la misma que la de los alrededores.
Equilibrio mecánico: la presión del sistema es la misma que la de los alrededores.


\subsection{Variables termodinámicas}

Las variables termodinámicas o variables de estado son las magnitudes que se emplean para describir el estado de un sistema termodinámico. Dependiendo de la naturaleza del sistema termodinámico objeto de estudio, pueden elegirse distintos conjuntos de variables termodinámicas para describirlo. En el caso de un gas, estas variables son:

Masa (m ó n): es la cantidad de sustancia que tiene el sistema. En el Sistema Internacional se expresa respectivamente en kilogramos (kg) o en número de moles (mol).
Volumen (V): es el espacio tridimensional que ocupa el sistema. En el Sistema Internacional se expresa en metros cúbicos (m3). Si bien el litro (l) no es una unidad del Sistema Internacional, es ampliamente utilizada. Su conversión a metros cúbicos es: 1 l = 10-3 m3.
Presión (p): Es la fuerza por unidad de área aplicada sobre un cuerpo en la dirección perpendicular a su superficie. En el Sistema Internacional se expresa en pascales (Pa). La atmósfera es una unidad de presión comúnmente utilizada. Su conversión a pascales es: $1 atm  \approx 105 Pa.$
Temperatura (T ó t): A nivel microscópico la temperatura de un sistema está relacionada con la energía cinética que tienen las moléculas que lo constituyen. Macroscópicamente, la temperatura es una magnitud que determina el sentido en que se produce el flujo de calor cuando dos cuerpos se ponen en contacto. En el Sistema Internacional se mide en kelvin (K), aunque la escala Celsius se emplea con frecuencia. La conversión entre las dos escalas es: T (K) = t (ºC) + 273.
En la siguiente figura se ha representado un gas encerrado en un recipiente y las variables termodinámicas que describen su estado.

\begin{center}
\includegraphics{images/02/variables_termo.png}    
\end{center}

Cuando un sistema se encuentra en equilibrio, las variables termodinámicas están relacionadas mediante una ecuación denominada ecuación de estado.

Variables extensivas e intensivas
En termodinámica, una variable extensiva es una magnitud cuyo valor es proporcional al tamaño del sistema que describe. Esta magnitud puede ser expresada como suma de las magnitudes de un conjunto de subsistemas que formen el sistema original. Por ejemplo la masa y el volumen son variables extensivas.

Una variable intensiva es aquella cuyo valor no depende del tamaño ni la cantidad de materia del sistema. Es decir, tiene el mismo valor para un sistema que para cada una de sus partes consideradas como subsistemas del mismo. La temperatura y la presión son variables intensivas.

Función de estado
Una función de estado es una propiedad de un sistema termodinámico que depende sólo del estado del sistema, y no de la forma en que el sistema llegó a dicho estado. Por ejemplo, la energía interna y la entropía son funciones de estado.

El calor y el trabajo no son funciones de estado, ya que su valor depende del tipo de transformación que experimenta un sistema desde su estado inicial a su estado final.

Las funciones de estado pueden verse como propiedades del sistema, mientras que las funciones que no son de estado representan procesos en los que las funciones de estado varían.
\end{document}